\documentclass[11pt,letterpaper]{article}

\usepackage[margin=1in]{geometry}
\usepackage[T1]{fontenc}
\usepackage{lmodern}
\usepackage{microtype}
\usepackage{amsmath,amssymb}
\usepackage{booktabs,tabularx,array}
\usepackage{xcolor}
\usepackage{enumitem}
\usepackage{hyperref}
\hypersetup{colorlinks=true,linkcolor=blue!45!black,urlcolor=blue!45!black,
  pdftitle={Reinforcement Learning Portfolio Management: Project Progress}}
\urlstyle{tt}
\setlength{\parindent}{0pt}
\setlength{\parskip}{6pt}
\setlength{\emergencystretch}{2em}
\setlist[itemize]{leftmargin=*,itemsep=3pt,topsep=3pt}
\newcolumntype{Y}{>{\raggedright\arraybackslash}X}

\title{Reinforcement Learning Portfolio Management\\[4pt]
  \large Project Progress Report}
\author{RL Portfolio Management Project}
\date{October 15, 2026}

\begin{document}
\maketitle

\section{Introduction}

\textcolor{red}{[THOMAS]}


This project investigates whether a reinforcement learning agent can improve
portfolio allocation across U.S. equity sectors. The agent observes historical
market signals and its current holdings, then adjusts sector weights every five
trading days. Portfolio returns, trading costs, borrowing costs, and weight drift
are calculated daily. The objective is to maximize portfolio growth after the
modeled implementation costs.

The current implementation is a research pipeline consisting of market-data collection,
feature engineering, a custom portfolio environment, PPO training, saved model
checkpoints, and development-validation comparisons against simple benchmarks.
Three main model stages have been implemented: a long-only allocator, an unscaled
long-short allocator, and a long-short allocator with scaled changes to portfolio
weights.

The downloaded market history ends on September
30, 2026. All performance figures below are historical development-validation
results used during model development as the ETF experiments reserve the
final two years for a later evaluation.

\section{Technology Used}

The project uses Python and Jupyter notebooks, with PyTorch supporting the neural
networks and Stable-Baselines3 providing the Proximal Policy Optimization (PPO)
implementation. Gymnasium supplies the interface for the custom allocation
environment. NumPy and pandas handle numerical operations, date alignment,
feature construction, and saved results. yfinance retrieves Yahoo Finance data.
Matplotlib supports the analysis plots.

\begin{table}[htbp]
\centering
\small
\begin{tabularx}{\linewidth}{@{}l l Y@{}}
\toprule
Component & Pinned version & Role \\
\midrule
Python & 3.11 (tested) & Runtime documented in the requirements file \\
NumPy / pandas & 1.26.4 / 2.2.3 & Arrays, feature panels, and performance analysis \\
yfinance & 0.2.61 & Adjusted prices and preliminary stock screening \\
Gymnasium & 0.29.1 & Observation, action, and episode interfaces \\
Stable-Baselines3 & 2.3.2 & PPO training and deterministic evaluation \\
PyTorch & 2.4.0 & Neural-network training \\
Matplotlib & 3.9.2 & Return, drawdown, and allocation plots \\
JupyterLab / ipykernel & 4.2.5 / 6.29.5 & Interactive experimentation \\
\bottomrule
\end{tabularx}
\caption{Dependencies recorded in \texttt{requirements.txt}.}
\end{table}

The workflow is organized into six numbered notebooks: data download, feature
engineering, agent training, benchmark analysis, action-scale search, and the
preliminary stock-fund experiment. The active settings and action/reward functions
are defined in \path{notebooks/experiment_config.py}. Saved experiment folders
contain checkpoints, configuration snapshots, chronological split summaries,
daily validation records, and aggregate metrics. 

\section{Data}

\subsection{Market history and investment universe}

The main dataset contains daily observations for the nine original sector ETFs, plus SPY and the VIX as context series. SPY and VIX provide information
to the policy but are not tradable assets in the current environment.

\begin{table}[htbp]
\centering
\small
\begin{tabular}{@{}ll@{}}
\toprule
ETF & Sector \\
\midrule
XLB & Materials \\
XLE & Energy \\
XLF & Financials \\
XLI & Industrials \\
XLK & Information Technology \\
XLP & Consumer Staples \\
XLU & Utilities \\
XLV & Health Care \\
XLY & Consumer Discretionary \\
\bottomrule
\end{tabular}
\caption{The current nine-asset allocation universe.}
\end{table}

To run the analysis, you must clone the GitHub repository and run the cells in \path{01_download_data.ipynb}
The download uses split- and dividend-adjusted prices
(\texttt{auto\_adjust=True}), with open, high, low, close, volume, and corporate-action
fields retained. The simulator uses adjusted close-to-close returns.

\subsection{Features and observations}

Let $P_{i,t}$ denote the adjusted closing price of asset $i$ on trading day $t$,
and let $N=9$ be the number of tradable sector ETFs. All windows below count
trading days: 5 days approximate one week, 21 days one month, and 63 days one
quarter. These horizons provide different views of recent market behavior.

\textcolor{red}{[ADD MORE?]}

\textbf{Features for each tradable ETF.} Each ETF contributes the following
eight features:

\begin{itemize}
  \item \textbf{5-day return (\texttt{return\_5d}).}
  \[
    r^{(5)}_{i,t}=\frac{P_{i,t}}{P_{i,t-5}}-1.
  \]
  This measures the price change over approximately one trading week. 

  \item \textbf{21-day return (\texttt{return\_21d}).}
  \[
    r^{(21)}_{i,t}=\frac{P_{i,t}}{P_{i,t-21}}-1.
  \]
  This measures approximately one month of performance. 

  \item \textbf{63-day return (\texttt{return\_63d}).}
  \[
    r^{(63)}_{i,t}=\frac{P_{i,t}}{P_{i,t-63}}-1.
  \]
  This measures approximately one quarter of performance. 

  \item \textbf{21-day annualized volatility
  (\texttt{volatility\_21d}).} 
  
  
  First define the daily log return
  $\ell_{i,t}=\log(P_{i,t}/P_{i,t-1})$. Then
  \[
    \overline\ell_{i,t}=\frac{1}{21}\sum_{k=0}^{20}\ell_{i,t-k},
    \qquad
    \sigma_{i,t}^{(21)}=
      \sqrt{252}\sqrt{\frac{1}{20}\sum_{k=0}^{20}
        (\ell_{i,t-k}-\overline\ell_{i,t})^2}.
  \]

  This measures recent return variability.
  The one-month window keeps the estimate responsive to recent conditions with 
  annualization putting it on a familiar yearly scale.

  \item \textbf{Drawdown relative to the 63-day high
  (\texttt{rolling\_drawdown\_63d}).}
  \[
    D_{i,t}^{(63)}=
      \frac{P_{i,t}}{\max_{0\leq k\leq62}P_{i,t-k}}-1.
  \]
  This describes how far a sector has fallen from its recent peak. The quarterly window
  gives context beyond the most recent month. This feature is the current
  distance from a rolling high.

  \item \textbf{Cross-sectional rank of the 21-day return
  (\texttt{cross\_sectional\_rank\_return\_21d}).}
  \[
    Q_{i,t}=\frac{
      \operatorname{rank}_{j=1,\ldots,N}(r^{(21)}_{j,t})_i}{N}.
  \]
  This identifies relative sector leadership even when all sectors rise or fall together. 

  \item \textbf{21-day return relative to SPY
  (\texttt{market\_relative\_return\_21d}).}
  \[
    M_{i,t}^{(21)}=r^{(21)}_{i,t}-r^{(21)}_{\mathrm{SPY},t}.
  \]
  This measures whether a sector has outperformed or underperformed the broad
  market over the month. The feature helps separate sector leadership from
  a market-wide rally.

  \item \textbf{63-day rolling market beta
  (\texttt{rolling\_market\_beta\_63d}).}
  \[
    \beta_{i,t}^{(63)}=
      \frac{\operatorname{Cov}_{63}
        (r^{(1)}_{i},r^{(1)}_{\mathrm{SPY}})_t}
      {\operatorname{Var}_{63}(r^{(1)}_{\mathrm{SPY}})_t},
    \qquad
    r^{(1)}_{i,t}=\frac{P_{i,t}}{P_{i,t-1}}-1.
  \]
  Covariance and variance use the latest 63 daily simple returns, including
  day $t$, with sample estimates. The quarterly window provides more observations
  for estimating this relationship than a one-month window.

\end{itemize}

\textbf{SPY context features.} SPY contributes its 5-, 21-, and 63-day returns,
21-day annualized volatility, and 63-day rolling drawdown. These five inputs describe broad-market
trend, variability, and distance from recent highs, helping the policy interpret
sector signals in their market context. 

\textbf{Market-wide features.} The following five features summarize the sector
universe and the volatility environment:

\begin{itemize}
  \item \textbf{Sector breadth above the 63-day moving average
  (\texttt{sector\_breadth\_above\_sma\_63d}).}
  \[
    \operatorname{SMA}_{i,t}^{(63)}=
      \frac{1}{63}\sum_{k=0}^{62}P_{i,t-k},
    \qquad
    B_t=\frac{1}{N}\sum_{i=1}^{N}
      \mathbf{1}\!\left\{P_{i,t}>\operatorname{SMA}_{i,t}^{(63)}\right\}.
  \]
  This is the fraction of sectors above their quarterly moving averages.
  Values near one indicate widespread participation in an upward trend with
  values near zero indicate broad weakness.

  \item \textbf{Dispersion of 21-day sector returns
  (\texttt{sector\_return\_dispersion\_21d}).}
  \[
    \overline r_t^{(21)}=\frac{1}{N}\sum_{i=1}^{N}r^{(21)}_{i,t},
    \qquad
    S_t=\sqrt{\frac{1}{N-1}\sum_{i=1}^{N}
      (r^{(21)}_{i,t}-\overline r_t^{(21)})^2}.
  \]
  This is the sample standard deviation across the nine sectors' monthly
  returns on the same date. It measures return differences across sectors, whereas each ETF's
  volatility measures variation through time.

  \item \textbf{Average pairwise 63-day sector correlation
  (\texttt{sector\_average\_correlation\_63d}).}
  \[
    C_t=\frac{2}{N(N-1)}\sum_{i<j}
      \operatorname{Corr}_{63}(r^{(1)}_i,r^{(1)}_j)_t.
  \]
  This describes how strongly sectors move together. Higher
  correlation suggests less diversification across sectors, while lower
  correlation indicates more independent movements.

  \item \textbf{VIX level (\texttt{vix\_level}).}
  \[
    V_t=\operatorname{Close}_{\texttt{\string^VIX},t}.
  \]
  This takes the downloaded VIX closing index level directly, without
  converting it into an asset return.

  \item \textbf{5-day VIX return (\texttt{vix\_return\_5d}).}
  \[
    v_t^{(5)}=\frac{V_t}{V_{t-5}}-1.
  \]
  This is the proportional change in the VIX index over one trading week.
  It is an index change, rather than the return from an investable VIX position.
\end{itemize}

\textbf{Portfolio accounting and observation preparation.}

\begin{itemize}
  \item \textbf{One-day simple return (\texttt{return\_1d}).}
  \[
    r^{(1)}_{i,t}=\frac{P_{i,t}}{P_{i,t-1}}-1.
  \]
  This is stored to calculate daily portfolio returns and weight drift, and
  it underlies beta and correlation estimates. 

  \item \textbf{Trailing standardization and clipping.} For each market-feature
  column $x_t$, set $W=1{,}260$ and compute
  \begin{align*}
    \mu_t &= \frac{1}{W}\sum_{k=1}^{W}x_{t-k}, \\
    s_t &= \sqrt{\frac{1}{W}\sum_{k=1}^{W}(x_{t-k}-\mu_t)^2}, \\
    z_t &= \min\!\left(10,\max\!\left(-10,
      \frac{x_t-\mu_t}{s_t}\right)\right).
  \end{align*}
  The scaling code uses a population standard deviation
  (\texttt{ddof=0}), unlike the sample standard deviations used in the
  volatility and dispersion features. Scaling puts inputs with different
  units, such as VIX levels and decimal returns, on comparable numerical
  scales. The five-year window provides a substantial historical reference
  that adapts over time. Statistics exclude the current observation and all
  future observations, and the same trailing procedure is used during
  training and evaluation. Clipping bounds extreme standardized values to
  support numerical stability; it does not remove their observations.

  \item \textbf{Current portfolio weights.} The nine current sector weights
  $w_{i,t^-}$ are appended directly to the scaled market features without
  trailing standardization. They show the policy its existing allocation,
  including any negative positions, so residual actions can depend on
  current holdings and the cost of changing them.
\end{itemize}

\subsection{Chronological partitions}

The most recent two years are reserved as a final holdout. Earlier observations
are split chronologically, with 80\% assigned to training before the scaling
warm-up is removed. Consequently, the usable training rows are fewer than 80\%
of the final, standardized development dataset.

\begin{table}[htbp]
\centering
\small
\begin{tabular}{@{}llll@{}}
\toprule
Partition & First date & Last date & Dates \\
\midrule
Training & 2004-04-08 & 2019-08-23 & 3,871 \\
Development validation & 2019-08-26 & 2024-09-30 & 1,283 \\
Reserved final holdout & 2024-10-01 & 2026-09-30 & 501 \\
\bottomrule
\end{tabular}
\caption{Partitions recorded in the saved model split files.}
\end{table}

Validation starts with an observation on August 26, 2019, so the saved return
series contains 1,282 realized daily returns from August 27, 2019 through September
30, 2024. The tables in this report use that common return period.

\section{Training}

\subsection{Proximal Policy Optimization (PPO)}

\textcolor{red}{[ETHAN]}

Proximal Policy Optimization (PPO) is a reinforcement learning algorithm that
learns a decision rule, called a \emph{policy}, through interaction with an
environment. In this project, the environment is the historical portfolio
simulator: the agent observes market features and current holdings, proposes
allocation changes, and receives a reward based on the resulting portfolio
growth after costs. Training seeks to improve expected discounted cumulative
reward across successive decisions.

PPO uses an \emph{actor--critic} approach:
\begin{itemize}
  \item \textbf{Actor (policy).} The policy $\pi_\theta(a\mid s)$ defines a
  distribution over actions given observation $s$, with trainable parameters
  $\theta$. Sampling actions during training provides exploration. Here, the
  sampled actions are converted into changes in the nine sector weights.

  \item \textbf{Critic (value function).} The value function $V_\phi(s)$ estimates
  the expected discounted future reward from observation $s$. It supplies a
  baseline for assessing actions, and is trained against estimated return targets.

  \item \textbf{Advantage estimate.} The estimate $\widehat A_k$ measures how
  favorable an action's outcome was relative to the critic's baseline, using
  subsequent rewards and value estimates. Positive advantages encourage the
  actor to make those actions more likely; negative advantages encourage the
  opposite. Here, $k$ indexes policy decisions, each normally spanning five
  trading days.
\end{itemize}

PPO collects a batch of experience with the current policy, then performs
multiple minibatch optimization passes before collecting a new batch. Its
clipped policy objective discourages excessive changes to action likelihoods
during those updates. Define the likelihood ratio
\textcolor{red}{[VERIFY THIS]}

\[
  q_k(\theta)=
    \frac{\pi_\theta(a_k\mid s_k)}
         {\pi_{\theta_{\mathrm{old}}}(a_k\mid s_k)}.
\]

\textcolor{red}{[VERIFY THIS]}
For continuous actions, the numerator and denominator are probability densities.
The actor maximizes
\[
  L^{\mathrm{clip}}(\theta)=\widehat{\mathbb E}_k
    \left[\min\!\left(
      q_k(\theta)\widehat A_k,\,
      \operatorname{clip}(q_k(\theta),1-\epsilon,1+\epsilon)
        \widehat A_k
    \right)\right],
\]
where $\epsilon$ controls clipping and $\widehat{\mathbb E}_k$ denotes the batch
average. The minimum removes the incentive to push an already favorable
likelihood change beyond the clipping threshold. This encourages conservative
updates but does not impose a hard bound on total policy change.\footnote{Schulman
et al. (2017), \emph{Proximal Policy Optimization Algorithms},
\url{https://arxiv.org/abs/1707.06347}.}

\textcolor{red}{[VERIFY THIS]}
PPO is an on-policy method: new experience is generated with the latest policy
after each update cycle. The Stable-Baselines3 implementation also normalizes
advantages during training.\footnote{Stable-Baselines3 2.3.2, PPO documentation:
\url{https://stable-baselines3.readthedocs.io/en/v2.3.2/modules/ppo.html}.}
We use it as a practical starting point because it supports continuous actions
and integrates with our Gymnasium simulator.


\subsection{Learning setup}

All models use PPO with a feedforward multilayer-perceptron policy
(\texttt{MlpPolicy}). The policy outputs nine continuous actions in $[-1,1]$.
The model-specific action function converts these actions into
changes in portfolio weights. Episodes begin with an equal-weight portfolio
 and sample random starting points within the training partition. Each 
 episode spans up to five trading years (1,260 daily intervals), while 
 the policy acts once every five trading days.


 \textcolor{red}{[HYPERPARAMETERS THAT COULD BE CHANGED]}

\begin{table}[htbp]
\centering
\small
\begin{tabularx}{\linewidth}{@{}l Y@{}}
\toprule
Setting & Current value \\
\midrule
Requested training budget & 300,000 environment steps per run \\
Learning rate & $3\times10^{-4}$ \\
PPO rollout length & 2,048 environment steps \\
Minibatch size & 64 \\
Episode length & Five trading years, with random training starts \\
Decision interval & Five trading days; daily valuation and accounting \\
Initial main-model seed & 42 \\
Repeated-seed experiments & 42, 7, 21, 123, and 2026 \\
Evaluation policy & Deterministic actions over development validation \\
\bottomrule
\end{tabularx}
\caption{Shared training settings. An environment step normally spans five
trading days; it is not a single daily return. PPO can round the requested
training budget up to finish a rollout.}
\end{table}

\subsection{Portfolio accounting and reward}

Let $w_t$ be the weights established for a daily return interval, $w_{t^-}$ the
weights before trading, and $r_{t+1}$ the vector of next-day asset returns.
Turnover is the full sum of absolute weight changes, counting both purchases
and sales:
\[
  \tau_t = \sum_{i=1}^{N} \left|w_{i,t} - w_{i,t^-}\right|
\]

 \textcolor{red}{[ASSUMING TRANSACTION COSTS AND BORROWING, COULD BE CHANGED]}

The environment charges a transaction-cost rate $c=0.0005$ (5 basis points per
dollar traded), plus a daily borrowing charge based on an annual rate $b=0.01$
applied to short exposure. Net return and the daily reward are
\begin{align*}
  r^{\mathrm{net}}_{t+1}
    &= w_t^\top r_{t+1}
       - c\tau_t
       - \frac{b}{252}\sum_i\max(-w_{i,t},0), \\
  R_{t+1} &= 100\log\!\left(1+r^{\mathrm{net}}_{t+1}\right).
\end{align*}
The PPO step reward sums these daily rewards across the decision interval.
The objective therefore rewards compounded growth after costs, with no separate
drawdown or volatility penalty. Between decisions, holdings drift with asset
returns rather than being reset to constant weights every day.

Performance reporting includes compounded net return, annualized net return,
annualized volatility, Sharpe ratio, maximum drawdown, turnover, exposure, and
implementation costs. Annualization uses 252 trading days. The Sharpe ratio uses
a zero risk-free rate. Reported annualized implementation cost is the mean daily
transaction-plus-borrow charge multiplied by 252.


\section{Model 1: Long Only}

 \textcolor{red}{[HANK]}

The first model learns changes to a fully invested long-only portfolio.
For action vector $a_t$, the proposed weights are
\[
  \widetilde w_t = w_{t^-} + a_t-\overline a_t\mathbf{1}.
\]
Subtracting the mean action makes the proposed change sum to zero. The model can 
concentrate in selected sectors, but it cannot short, borrow to increase exposure, or allocate to a cash asset.

The saved seed-42 run in \path{models/long_only/} achieved a 23.51\% annualized
net return, a 0.975 Sharpe ratio, and a maximum drawdown of 35.07\%.
Its annualized turnover was 22.54 times portfolio value, with an annualized
implementation cost of 1.13\%.

\section{Model 2: Long Short}

The second model retains the same centered, unscaled residual action but permits
negative weights which allows the agent to take short position. Net exposure remains approximately 100\%. If a proposed
portfolio exceeds 160\% gross exposure, its positive positions are rescaled to
130\% long exposure and its negative positions to 30\% short exposure. When the
proposal is already within the gross limit, it is retained. Thus the model can
use less than the full short allowance.

The seed-42 run in \path{models/long_short/} produced a 19.37\% annualized net
return and a Sharpe ratio of 0.932, with a 33.10\% maximum drawdown.
Annualized turnover increased to 45.74 times portfolio value, and implementation
cost rose to 2.59\% per year. Shorting expanded the available portfolios, but the
large residual actions made trading expensive, and this run did not improve on
the long-only run's return or Sharpe ratio.

The exposure constraint is applied when actions establish new target weights.
Daily drift can subsequently push gross exposure above 160\%; the original
long-short run records a maximum of approximately 164.71\%. The current simulator
therefore does not enforce a continuous daily gross-exposure ceiling.

\section{Model 3: Long Short with Scaled Actions}

\subsection{Smaller changes to current holdings}

The third model introduces an action scale $\alpha$:
\[
  \widetilde w_t
    = w_{t^-} + \alpha\left(a_t-\overline a_t\mathbf{1}\right).
\]
It uses the same long-short exposure mapping as Model 2. Scaling changes how far
the policy moves from its current holdings at each decision; it does not impose
a fixed final position size or a hard turnover cap.

The scaled run used $\alpha=0.10$. For seed 42, annualized turnover fell
from 45.74 to 16.69, a reduction of approximately 63.5\% relative to the
unscaled long-short run. Annualized implementation cost fell from 2.59\%
to 1.13\%. Annualized net return increased to 21.97\%, and Sharpe increased
slightly to 0.936. However, maximum drawdown worsened from 33.10\% to
34.74\% in magnitude.

\begin{table}[htbp]
\centering
\small
\setlength{\tabcolsep}{4pt}
\begin{tabular}{@{}lrrr@{}}
\toprule
Metric & Long only & Long short & Scaled 0.10 \\
\midrule
Cumulative net return & 192.74\% & 146.12\% & 174.67\% \\
Annualized net return & 23.51\% & 19.37\% & 21.97\% \\
Annualized volatility & 24.83\% & 21.47\% & 24.41\% \\
Sharpe ratio & 0.975 & 0.932 & 0.936 \\
Maximum drawdown & $-35.07$\% & $-33.10$\% & $-34.74$\% \\
Annualized turnover & 22.54 & 45.74 & 16.69 \\
Annualized impl.\ cost & 1.13\% & 2.59\% & 1.13\% \\
\bottomrule
\end{tabular}
\caption{Saved seed-42 development-validation results. The scaled model uses
$\alpha=0.10$. Turnover is expressed as multiples of portfolio value per year.}
\end{table}


\end{document}
